\par\Needspace{6\baselineskip}\medskip
\begin{center}\hypertarget{native-231-article168}{}\pdfbookmark[2]{168.}{article-168}168.\end{center}
Combinando aequationes 1 et 2 art. 112, prodeunt in signis disquisitionis praesentis aequationes sequentes:
\[
\begin{aligned}
0={}&(n12)R\cos B\sin(l-a)\\
&-(n02)\bigl(\rho'\cos\beta'\sin(a'-a)+R'\cos B'\sin(l'-a)\bigr)\\
&+(n01)\bigl(\rho''\cos\beta''\sin(a''-a)+R''\cos B''\sin(l''-a)\bigr)
\end{aligned}
\]
\SourceLeaf{232}
\[
\begin{aligned}
0={}&(n23)\bigl(\rho'\cos\beta'\sin(a'''-a')+R'\cos B'\sin(a'''-l')\bigr)\\
&-(n13)\bigl(\rho''\cos\beta''\sin(a'''-a'')+R''\cos B''\sin(a'''-l'')\bigr)\\
&+(n12)R'''\cos B'''\sin(a'''-l''').
\end{aligned}
\]

Hae aequationes, statuendo
\[
\begin{gathered}
\frac{R'\cos B'\sin(l'-a)}{\cos\beta'\sin(a'-a)}-R'\cos\delta'=b'\\[.7em]
\frac{R\cos B\sin(l-a)}{\cos\beta''\sin(a''-a)}=\lambda\\[.7em]
\frac{R''\cos B''\sin(a'''-l'')}{\cos\beta''\sin(a'''-a'')}-R''\cos\delta''=b''\\[.7em]
\frac{R'''\cos B'''\sin(a'''-l''')}{\cos\beta'\sin(a'''-a')}=\lambda'''\\[.7em]
\frac{R'\cos B'\sin(a'''-l')}{\cos\beta'\sin(a'''-a')}-R'\cos\delta'=\varkappa'\\[.7em]
\frac{\cos\beta'\sin(a'-a)}{\cos\beta''\sin(a''-a)}=\mu'\\[.7em]
\frac{R''\cos B''\sin(l''-a)}{\cos\beta''\sin(a''-a)}-R''\cos\delta''=\varkappa''\\[.7em]
\frac{\cos\beta''\sin(a'''-a'')}{\cos\beta'\sin(a'''-a')}=\mu''
\end{gathered}
\]
omnibusque rite reductis, transeunt in sequentes
\[
\begin{gathered}
\frac{\mu'(1+P')(x'+b')}{1+\dfrac{Q'}{(x'x'+a'a')^{\frac32}}}=x''+\varkappa''+\lambda P'\\[1em]
\frac{\mu''(1+P'')(x''+b'')}{1+\dfrac{Q''}{(x''x''+a''a'')^{\frac32}}}=x'+\varkappa'+\lambda'''P'',
\end{gathered}
\]
sive, statuendo insuper
\[
\begin{gathered}
-\varkappa''-\lambda P'=c',\qquad\mu'(1+P')=d'\\[.5em]
-\varkappa'-\lambda'''P''=c'',\qquad\mu''(1+P'')=d'',
\end{gathered}
\]
in hasce
\[
\tag*{I.}x''=c'+\frac{d'(x'+b')}{1+\dfrac{Q'}{(x'x'+a'a')^{\frac32}}}
\]
\[
\tag*{II.}x'=c''+\frac{d''(x''+b'')}{1+\dfrac{Q''}{(x''x''+a''a'')^{\frac32}}}\mathpunct{\cdot}
\]
Adiumento harum duarum aequationum $x'$ et $x''$ ex \newline $a',b',c',d',Q',a'',b'',c'',d'',Q''$ determinari poterunt. Quodsi quidem $x'$ vel $x''$ inde eliminanda esset, ad aequationem ordinis permagni\HistoricalNoteEight{Handschriftliche Bemerkung:] ordinis $64^{\mathrm{ti}}$.} delaberemur: attamen per methodos indirectas incognitarum $x',x''$ valores ex illis aequationibus forma non mutata \SourceLeaf{233}satis expedite elicientur. Plerumque valores incognitarum approximati iam prodeunt, si primo $Q'$ atque $Q''$ negliguntur; scilicet
\[
\begin{gathered}
x'=\frac{c''+d''(b''+c')+d'd''b'}{1-d'd''}\\[.6em]
x''=\frac{c'+d'(b'+c'')+d'd''b''}{1-d'd''}\mathpunct{\cdot}
\end{gathered}
\]
Quamprimum autem valor approximatus alterutrius incognitae habetur, valores aequationibus exacte satisfacientes facillime elicientur. Sit scilicet $\xi'$ valor approximatus ipsius $x'$, quo in aequatione I substituto prodeat $x''=\xi''$; perinde substituto $x''=\xi''$ in aequatione II prodeat inde $x'=X'$; repetantur eaedem operationes, substituendo pro $x'$ in I valorem alium $\xi'+\nu'$, unde prodeat $x''=\xi''+\nu''$, quo valore in II substituto prodeat inde $x'=X'+N'$. Tum valor correctus ipsius $x'$ erit
\[
=\xi'+\frac{(\xi'-X')\nu'}{N'-\nu'}=\frac{\xi'N'-X'\nu'}{N'-\nu'}\mathpunct{\raisebox{.45ex}{,}}
\]
valorque correctus ipsius $x''$
\[
=\xi''+\frac{(\xi'-X')\nu''}{N'-\nu'}\mathpunct{\cdot}
\]
Si operae pretium videtur, cum valore correcto ipsius $x'$ alioque levius mutato eaedem operationes repetentur, donec valores ipsarum $x',x''$ aequationibus I, II exacte satisfacientes prodierint. Ceterum analystae vel mediocriter tantum exercitato subsidia calculum contrahendi haud deerunt.

In his operationibus quantitates irrationales $(x'x'+a'a')^{\frac32}$, $(x''x''+a''a'')^{\frac32}$ commode calculantur per introductionem arcuum $z',z''$, quorum tangentes resp. sunt $\dfrac{a'}{x'}\mathpunct{\raisebox{.45ex}{,}}$ $\dfrac{a''}{x''}\mathpunct{\raisebox{.45ex}{,}}$ unde fit
\[
\begin{gathered}
\surd(x'x'+a'a'\mathpunct{\raisebox{-.5ex}{.}})=r'=\frac{a'}{\sin z'}=\frac{x'}{\cos z'}\\[.6em]
\surd(x''x''+a''a'')=r''=\frac{a''}{\sin z''}=\frac{x''}{\cos z''}\mathpunct{\cdot}
\end{gathered}
\]
Hi arcus auxiliares, quos inter $0$ et $180^0$ accipere oportet, ut $r',r''$ positivi evadant, manifesto cum arcubus $C'B'$, $C''B''$ identici erunt, unde patet, hacce ratione non modo $r'$ et $r''$, sed etiam situm punctorum $C',C''$ innotescere.

Haecce determinatio quantitatum $x',x''$ requirit, ut\newline $a',a'',b',b'',c',c'',d',d'',Q',Q''$ cognitae sint, quarum quantitatum quatuor primae quidem per problematis data habentur, quatuor sequentes autem a $P',P''$ pendent. Iam \SourceLeaf{234}quantitates $P',P'',Q',Q''$ exacte quidem nondum determinari possunt; attamen quum habeatur
\[
\tag*{III.}P'=\frac{t''-t'}{t'-t}\cdot\frac{(\eta01)}{(\eta12)}
\]
\[
\tag*{IV.}P''=\frac{t''-t'}{t'''-t''}\cdot\frac{(\eta23)}{(\eta12)}
\]
\[
\tag*{V.}\begin{aligned}
Q'={}&\tfrac12kk(t'-t)(t''-t')\cdot\frac{r'r'}{rr''}\cdot{}\\[.5em]
&\frac{1}{(\eta01)(\eta12)\cos\tfrac12(v'-v)\cos\tfrac12(v''-v)\cos\tfrac12(v''-v')}
\end{aligned}
\]
\[
\tag*{VI.}\begin{aligned}
Q''={}&\tfrac12kk(t''-t')(t'''-t'')\cdot\frac{r''r''}{r'r'''}\cdot{}\\[.5em]
&\frac{1}{(\eta12)(\eta23)\cos\tfrac12(v''-v')\cos\tfrac12(v'''-v')\cos\tfrac12(v'''-v'')}\mathpunct{\raisebox{.45ex}{,}}
\end{aligned}
\]
statim adsunt valores approximati
\[
\begin{gathered}
P'=\frac{t''-t'}{t'-t}\mathpunct{\raisebox{.45ex}{,}}\qquad P''=\frac{t''-t'}{t'''-t''}\\[.6em]
Q'=\tfrac12kk(t'-t)(t''-t'),\qquad Q''=\tfrac12kk(t''-t')(t'''-t''),
\end{gathered}
\]
quibus calculus primus superstruetur.
